\documentclass[5pt, a4paper]{article}

\usepackage[utf8]{inputenc}
\usepackage[T1]{fontenc}
\usepackage{textcomp}
\usepackage{amsmath, amssymb}
\usepackage{multicol}
\usepackage{proof}
\usepackage{enumitem}
\usepackage{fancyhdr}
\usepackage[margin=1in]{geometry}
% figure support
\usepackage{import}
\usepackage{xifthen}
\pdfminorversion=7

\usepackage[most]{tcolorbox}
\usepackage{pdfpages}
\usepackage{transparent}
\newcommand{\incfig}[1]{%
    \def\svgwidth{\columnwidth}
    \import{./figures/}{#1.pdf_tex}
}

\tcbset{
colback=gray!5,
colframe=blue!50!white,
coltitle=black,
arc=0pt,
boxrule=-0.1mm,
fonttitle=\bfseries,
coltitle=black,
left=2.5mm,
leftrule=1mm,
right=3.5mm,
colbacktitle=black!10!white,
toptitle=0.75mm,
bottomtitle=0.5mm,
}
\begin{document}
\section*{Chapter 5: Interpolation}
We'll start out with polynomial interpolation, which is basically try to fit a
polynomial such that it matches the given data points (we'l skip over trig and
rational interpolations in this course.)

In general, for a polynomial of degree $d$, a solution exists when the number of
the given data points is no more than $d + 1$, and the solution is unique as
long as there is exactly $d+1$ points. These all depends on the basis matrix.
One can choose any basis, each with their own advantages and disadvatages. I'll
start with the monomial basis matrix, also known as the Vandermonde matrix. It
looks something like this: \[
X = \begin{pmatrix} 
    1 & x_0 & x_0^2 & \cdots & x_0^d \\
    1 & x_1 & x_1^2 & \cdots & x_1^d \\
    \vdots & \vdots & \vdots & \ddots & \vdots \\
    1 & x_n & x_n^2 & \cdots & x_n^d \\
\end{pmatrix} 
= \begin{pmatrix} 
    M_0(x_0) & M_1(x_0) & M_2(x_0) & \cdots & M_d(x_0) \\
    M_0(x_1) & M_1(x_1) & M_2(x_1) & \cdots & M_d(x_1) \\
    \vdots & \vdots & \vdots & \ddots & \vdots \\
    M_0(x_n) & M_1(x_n) & M_2(x_n) & \cdots & M_d(x_n) \\
\end{pmatrix} 
.\] 
This looks like 

\begin{align*}
    p(x) &= a_0 + a_1x^2+ \ldots + a_n x^d \\
    &=  X \overline{a}
.\end{align*}

\paragraph{Example}
Find the polynomial that fit the points $(-2,-27),(0,-1),(1,0)$.
\[
\begin{pmatrix} 
    1 & x_0 & x_0^2 \\
    1 & x_1 & x_1^2 \\
    1 & x_2 & x_2^2 \\
\end{pmatrix} 
\begin{pmatrix} a_1 \\ a_2 \\ a_3 \end{pmatrix} 
=\begin{pmatrix} y_1 \\ y_2 \\ y_3 \end{pmatrix} 
.\] 
\[
\begin{pmatrix} 
    1 & -2 & 4  \\
    1 & 0 & 0 \\
    1 & 1 & 1 \\
\end{pmatrix} 
\begin{pmatrix} a_1 \\ a_2 \\ a_3 \end{pmatrix} 
=
\begin{pmatrix} -27 \\ -1 \\ 0 \end{pmatrix} 
.\] 

Let's consider another basis: the Langrange. S





Consider another basis: the Newton.
\end{document}
